\chapter{기계는 어떻게 움직임을 배우는가}
% full version: chapters/15_motorlearning.tex

지금까지 배우는 방법 두 가지를 보았다. 교사 없이 “함께 일하면 이어진다”는 규칙으로 배우면, 기계는 자주 일어나는 것을 기억한다. 도파민이 있으면 기계는 뜻밖의 일에서 배운다. 결과가 예상보다 좋았는가, 나빴는가. 그런데 농구 골대에 공을 던져 빗나갔을 때, 여러분은 결과가 나빴다는 것만 아는 게 아니다. \emph{얼마나}, \emph{어느 쪽으로} 빗나갔는지도 안다. 이것이 세 번째 교사다. 방향이 있는 오차, 그것도 출력마다 따로 주어지는 오차다.

\section*{오차 전선}

뇌에는 움직임 학습을 위한 별도의 회로가 있고, 그 구조는 아주 특이하다. 작은 \nt{GLUT} 뉴런 수백억 개가 입력 신호를 아주 긴 특징 목록으로 펼친다. 이 회로의 출력 뉴런은 억제성 \nt{GABA} 뉴런이며, 수천만 개가 있고 하나하나가 입력 10만 개를 듣는다. 그리고 출력 뉴런마다 특별한 전선이 딱 하나씩 닿는다. 바로 오차 전선이다. 어떤 입력이 활성인 동안 오차가 도착하면 그 입력은 약해진다. 엔지니어가 적응 필터를 학습시킬 때 쓰는 고전적인 “오차를 줄여라” 규칙이다.

\pencilfig{15}{\pencilart{15}}{빗나감은 “나쁘다”만이 아니라 “얼마나, 어느 쪽으로”도 알려 준다.}

이 교사는 도파민 교사보다 빠르다. 출력마다 자기 오차가 있으므로 모두에게 가는 공통 신호에서 자기 몫을 가려낼 필요가 없다. 다만 오차는 늦게 도착하므로, 연결에는 다시 자기가 참여했음을 기억하는 “스티커”가 필요하다. 실험에 따르면 이 스티커의 지속 시간은 오차의 지연에 맞춰져 있다.

\section*{내부 모델}

이런 회로는 무엇을 배울까? 자기 몸의 모델, 곧 어떤 명령이 어떤 결과를 낼지다. 이 모델이 있으면 늦게 오는 감지기를 기다릴 필요가 없다. 결과를 미리 예측할 수 있다. 한 가설에 따르면 스스로 간지럼을 태울 수 없는 까닭도 이것이다. 뇌는 무슨 일이 일어날지 미리 알고 예측한 것을 빼 버린다.

\section*{빠른 학습자와 느린 학습자}

그림을 옆으로 밀어 보이게 하는 안경을 쓰고 다트를 던져 보자. 처음 몇 번은 옆으로 빗나가지만, 수십 번 시도하면 적응한다. 안경을 벗으면 다시 빗나간다. 이번에는 반대쪽이다. 흥미로운 일은 그다음에 일어난다. 실험 결과는 학습자 둘이 동시에 있다고 보면 잘 설명된다. 빠른 학습자는 빨리 배우고 빨리 잊는다. 느린 학습자는 천천히 배우지만 오래 기억한다. 오랫동안 적응한 뒤 반대 방향으로 짧게 다시 배우면 전체 보정은 0으로 돌아간다. 그런데 그 뒤 옛 솜씨가 저절로 일부 되살아난다. 빠른 학습자는 잊었지만 느린 학습자는 기억하기 때문이다. 학습자가 하나인 모델은 이렇게 하지 못하지만 사람은 한다.

\pencilfig{115}{%
% figures/tworate_*.dat from models/motor_learning.py: adapt 200 throws, reverse until the sum is back to zero, then no errors at all
\begin{scope}[x=0.026cm, y=1.45cm]
\fill[pcGray, opacity=0.08] (220,-1.1) rectangle (226,1.1);
\draw[pen] (0,0) -- (340,0);
\draw[pcGray, line width=0.4pt] plot file {figures/tworate_p.dat};
\draw[penlite=pcAmber] plot file {figures/tworate_fast.dat};
\draw[penlite=pcBlue] plot file {figures/tworate_slow.dat};
\draw[pcGray!40!black, line width=1.1pt] plot file {figures/tworate_net.dat};
\draw[pcGray, densely dashed, line width=0.4pt] (226,-1.1) -- (226,1.1);
\end{scope}
\node[lbl, anchor=south] at (3.1,1.47) {안경 씀};
\node[lbl, anchor=north west, align=left] at (6.3,-0.55) {짧게\\재학습};
\node[lbl, anchor=south west] at (6.0,1.47) {오차 없음};
\node[lbl, text=pcBlue, anchor=north] at (4.4,0.8) {느린 쪽};
\node[lbl, text=pcAmber!80!black, anchor=south west] at (1.15,0.5) {빠른 쪽};
\node[lbl, text=pcGray!40!black, anchor=south] at (7.7,0.95) {합이 되살아남};
}{모델에서 두 학습자의 보정. 짧게 다시 배우고 나면 합은 0이지만, 느린 학습자가 옛것을 기억하므로 그것이 저절로 되살아난다.}

왜 학습자가 둘일까? 가설에 따르면 9장의 항해사가 쓰던 것과 같은 처방이다. 세상은 빠르게도 느리게도 변하므로 둘 다 지켜보는 것이 현명하다.

\fullversion{15}
